\documentclass[12pt]{article}

\usepackage[a4paper, margin=2.5cm]{geometry}
\usepackage{amsmath}
\usepackage{amssymb}
\usepackage{amsthm}
\usepackage[colorlinks=true, linkcolor=blue, citecolor=blue, urlcolor=blue]{hyperref}
\usepackage{booktabs}
\usepackage{natbib}

\newtheorem{theorem}{Theorem}
\newtheorem{proposition}[theorem]{Proposition}
\newtheorem{definition}[theorem]{Definition}
\newtheorem{remark}[theorem]{Remark}

\title{From Dirac Solutions to Physical Reality:\\Crystallisation as the Selection Mechanism}
\author{%
  Lee Smart\\[4pt]
  \textit{Institute of Vibrational Field Dynamics}\\[2pt]
  \texttt{contact@vibrationalfielddynamics.org}\\[2pt]
  \texttt{@vfd\_org}%
}
\date{April 2026}

\begin{document}

\maketitle

\noindent\textbf{Status:} Preprint (not peer-reviewed)

\begin{abstract}
The Dirac equation defines the space of admissible relativistic quantum states --- including particle and antiparticle branches, spin degrees of freedom, and momentum modes --- but does not specify the mechanism by which a single physical state is realised from this space. Standard quantum mechanics supplies outcome probabilities via the Born rule, but in orthodox formulations does not specify a universally accepted physical selection mechanism.

We propose \textit{crystallisation} as a deterministic, constraint-based selection architecture acting over a restricted admissible subset of the Dirac solution space. Crystallisation is defined as asymptotic convergence under a projected gradient flow toward a state that extremises a schematic constraint functional $F(\psi)$. The selected state is determined by the geometry of the constraint landscape, not by stochastic projection.

Within the intended interpretation, candidate stable particles correspond to local extrema of $F$; antiparticles are symmetry-related extrema on conjugate solution branches; and annihilation is provisionally interpreted, in a field-theoretic extension, as a transition toward a more stable admissible configuration. The proposal extends rather than replaces the Dirac equation and standard quantum mechanics. The constraint functional is schematic and not uniquely derived; the paper introduces the \textit{selection architecture}, not a claimed final covariant functional.
\end{abstract}

%==================================================================
\section{Introduction}\label{sec:intro}
%==================================================================

The Dirac equation~\citep{Dirac1928} is a cornerstone of relativistic quantum mechanics. Its successes include the natural emergence of electron spin, the prediction of antiparticles~\citep{Anderson1933}, and the correct fine-structure spectrum of hydrogen. Combined with quantum field theory, it provides the foundation for the Standard Model's treatment of fermionic matter.

The equation defines a precise mathematical space of admissible states: solutions to the relativistic wave equation satisfying specific boundary conditions, symmetry requirements, and normalisability constraints. This solution space contains particle and antiparticle branches, multiple spin orientations, and a continuum of momentum modes. All are admissible at the level of the equation, subject to the physical constraints imposed by normalisation, boundary conditions, and interpretation.

Yet the Dirac equation, like the Schr\"odinger equation before it, does not answer a fundamental question: \textit{given the space of admissible states, what determines which state is physically realised?}

Standard quantum mechanics supplies outcome probabilities via the Born rule: $P(\psi_i) = |\langle\psi_i|\psi\rangle|^2$. This assigns probabilities to outcomes but in orthodox formulations does not specify a universally accepted physical mechanism for the transition from superposition to definite outcome --- the quantum measurement problem~\citep{Schlosshauer2005}.

The present paper proposes a candidate resolution: a deterministic selection architecture, termed \textit{crystallisation}, that acts over a restricted admissible subset of the Dirac solution space. The architecture is defined by a schematic constraint functional whose extrema correspond to candidate physically realised states. The proposal extends rather than replaces the Dirac framework; it adds a selection layer to the existing mathematical structure without modifying the equation itself.

\subsection*{Scope and epistemic status}

This paper is a bridge paper: it introduces the formal architecture and physical interpretation of crystallisation in the context of relativistic quantum mechanics. It does not derive the complete functional form from first principles.

The paper works primarily at the level of a \textit{Dirac-inspired selection framework}. Interpretations involving annihilation and particle creation are heuristic extensions to field-theoretic settings; a full QFT formulation is deferred. The constraint functional is schematic and proposed, not uniquely determined. The paper should be read as introducing a \textit{selection architecture} with testable structural implications, not as a finished theory.

%==================================================================
\section{The Dirac Solution Space}\label{sec:dirac}
%==================================================================

\subsection{The equation}

The Dirac equation for a free fermion of mass $m$ is:
\begin{equation}\label{eq:dirac}
    (i\gamma^\mu \partial_\mu - m)\psi = 0
\end{equation}
where $\gamma^\mu$ are the Dirac gamma matrices satisfying $\{\gamma^\mu, \gamma^\nu\} = 2g^{\mu\nu}$, and $\psi$ is a four-component spinor field.

\subsection{The solution space}

The set of all solutions forms the Dirac solution space:
\begin{equation}\label{eq:solspace}
    \mathcal{S} = \left\{ \psi \;\middle|\; (i\gamma^\mu \partial_\mu - m)\psi = 0 \right\}
\end{equation}

This space has the following structure:
\begin{itemize}
    \item \textbf{Particle branch:} positive-energy solutions $\psi^{(+)}$ with $E > 0$.
    \item \textbf{Antiparticle branch:} negative-energy solutions $\psi^{(-)}$ (reinterpreted via charge conjugation as positive-energy antiparticles).
    \item \textbf{Spin degrees of freedom:} each energy-momentum mode carries two spin orientations.
    \item \textbf{Momentum continuum:} solutions exist for all momenta $\mathbf{p}$ satisfying $E^2 = |\mathbf{p}|^2 + m^2$.
\end{itemize}

The equation constrains admissible states but does not specify a unique realised state.

\begin{remark}
In quantum field theory, the solution space is promoted to an operator algebra, and states are vectors in Fock space. The selection problem persists at the level of measurement outcomes. The present paper does not attempt a full Fock-space formulation; the field-theoretic extensions discussed in Section~\ref{sec:physics} are heuristic.
\end{remark}

%==================================================================
\section{The Missing Selection Mechanism}\label{sec:gap}
%==================================================================

\subsection{The selection problem}

Standard quantum mechanics provides the Born rule for outcome probabilities:
\begin{equation}
    P(\psi_i) = |\langle\psi_i|\psi\rangle|^2
\end{equation}

This assigns a probability distribution over outcomes, but no physical operator performs the selection:
\begin{equation}
    O_{\text{select}} : \mathcal{S} \to \psi^* \quad \text{(not provided by the orthodox formalism)}
\end{equation}

\begin{definition}[Selection Problem]
Given the solution space $\mathcal{S}$ of Eq.~\eqref{eq:solspace}, the selection problem asks: what determines the physically realised state $\psi^*$?
\end{definition}

\subsection{Existing approaches}

Several interpretations address this gap:
\begin{itemize}
    \item \textbf{Copenhagen:} provides operational closure through axiomatic collapse; does not specify an underlying mechanism~\citep{vonNeumann1932}.
    \item \textbf{Many-worlds:} avoids single-outcome collapse by positing universal branching; selection is reinterpreted as observer-relative~\citep{Everett1957}.
    \item \textbf{GRW / objective collapse:} introduces explicit stochastic modifications to the dynamics~\citep{GRW1986}.
    \item \textbf{Pilot wave:} provides a deterministic guidance equation with additional ontology (the pilot wave itself)~\citep{Bohm1952}.
\end{itemize}

None of these approaches formulates state selection as optimisation over a constraint functional on the Dirac solution space in the sense proposed here.

%==================================================================
\section{Crystallisation: Definition}\label{sec:crystallisation}
%==================================================================

\subsection{Admissible state domain}

The crystallisation mechanism does not act over the full solution space $\mathcal{S}$, which is infinite-dimensional and contains gauge-equivalent and non-normalisable configurations. Instead, it acts over a restricted admissible domain:

\begin{definition}[Admissible domain]\label{def:admissible}
The admissible domain $\mathcal{S}_{\mathrm{adm}} \subset \mathcal{S}$ consists of normalised, finite-energy, measurement-conditioned states satisfying the relevant boundary conditions and symmetry constraints of the physical context.
\end{definition}

The precise specification of $\mathcal{S}_{\mathrm{adm}}$ depends on the physical setting (e.g.\ scattering, bound state, measurement apparatus). The present paper does not provide a universal construction of $\mathcal{S}_{\mathrm{adm}}$; it assumes such a domain can be defined for each physical context.

\subsection{The schematic constraint functional}

The constraint functional $F$ encodes the physical requirements for a realised state. At the present level of development, $F$ is a \textit{schematic template} --- a structural proposal whose specific form is not uniquely derived:
\begin{equation}\label{eq:functional}
    F(\psi) = a_1\, C(\psi) - a_2\, H(\psi) + a_3\, R(\psi) - a_4\, V(\psi)
\end{equation}

where:
\begin{itemize}
    \item $C(\psi)$: \textbf{coherence measure} --- quantifies internal phase alignment and structural consistency.
    \item $H(\psi)$: \textbf{entropy / disorder} --- measures the degree of internal decoherence.
    \item $R(\psi)$: \textbf{resonance alignment} --- measures alignment with the constraint geometry.
    \item $V(\psi)$: \textbf{instability / variance} --- penalises fluctuations and structural fragility.
\end{itemize}

The coefficients $a_1, a_2, a_3, a_4 > 0$ are structural parameters. Their values, the specific operator definitions of $C$, $H$, $R$, $V$, and the Lorentz-covariance properties of $F$ are not established in the present paper.

\textsc{Status:} $F$ is a schematic functional template. The present paper introduces only the selection architecture, not the final covariant functional.

\subsection{Dynamical realisation}

The static selection statement $\psi^* \in \arg\sup_{\psi \in \mathcal{S}_{\mathrm{adm}}} F(\psi)$ is shorthand for asymptotic convergence under a projected gradient flow:
\begin{equation}\label{eq:flow}
    \frac{d\psi}{dt} = \Pi_{\mathcal{S}_{\mathrm{adm}}}\, \nabla_\psi F(\psi)
\end{equation}
where $\Pi_{\mathcal{S}_{\mathrm{adm}}}$ denotes projection onto the tangent space of $\mathcal{S}_{\mathrm{adm}}$ (ensuring the flow remains within the admissible domain).

\begin{definition}[Crystallisation]\label{def:crystal}
The crystallisation of an initial state $\psi_0 \in \mathcal{S}_{\mathrm{adm}}$ is defined as:
\begin{equation}\label{eq:crystal}
    \psi^* = \lim_{t \to \infty} \Phi_t(\psi_0)
\end{equation}
where $\Phi_t$ is the flow generated by Eq.~\eqref{eq:flow}, provided convergence exists.
\end{definition}

Whether this flow converges, whether the limit is unique, and whether the limit depends on $\psi_0$ (basin structure) are mathematical questions that depend on the specific form of $F$ and the topology of $\mathcal{S}_{\mathrm{adm}}$. The present paper does not resolve these questions; it defines the architecture and identifies the structural requirements.

\textsc{Status:} The crystallisation architecture is a formal definition. Existence and uniqueness of the flow limit, and the basin structure of $F$, are open mathematical problems within this framework.

%==================================================================
\section{Physical Interpretation}\label{sec:physics}
%==================================================================

\subsection{Particle stability}

Within the crystallisation framework, candidate stable particles correspond to local maxima of $F$ over $\mathcal{S}_{\mathrm{adm}}$. A state $\psi^*$ is a candidate stable configuration if:
\begin{equation}
    \frac{\delta^2 F}{\delta\psi^2}\bigg|_{\psi^*} < 0 \quad \text{(negative definite)}
\end{equation}

In the intended interpretation, distinct attractor basins of the flow~\eqref{eq:flow} would correspond to distinct stable realised state classes, potentially including particle types. The discreteness of the observed particle spectrum would then correspond to the discrete set of local maxima.

\subsection{Antiparticles}

The Dirac solution space $\mathcal{S}$ contains both particle ($E > 0$) and antiparticle ($E < 0$, reinterpreted) branches. In the crystallisation framework, antiparticles correspond to symmetry-related extrema of $F$ on the conjugate solution branch:
\begin{equation}
    \psi^*_{\text{particle}} \in \arg\sup_{\psi \in \mathcal{S}^{(+)}_{\mathrm{adm}}} F(\psi), \qquad
    \psi^*_{\text{anti}} \in \arg\sup_{\psi \in \mathcal{S}^{(-)}_{\mathrm{adm}}} F(\psi)
\end{equation}

If the effective constraint functional respects CPT-compatible symmetries, particle and antiparticle extrema may be related by the corresponding conjugation map: $F(\psi) = F(\mathcal{C}_{\text{CPT}}\psi)$.

\subsection{Annihilation (heuristic field-theoretic extension)}

Within a field-theoretic extension of this framework (not fully developed here), annihilation would be represented as:
\begin{equation}
    (\psi_1, \psi_2) \to \psi_{\text{lower}}
\end{equation}
where $\psi_1$ and $\psi_2$ are a particle-antiparticle pair and $\psi_{\text{lower}}$ is a configuration of higher global stability (typically radiation modes). This is a heuristic interpretation; the full QFT formulation is deferred to future work.

%==================================================================
\section{Relation to Standard Quantum Mechanics}\label{sec:relation}
%==================================================================

The crystallisation proposal extends rather than replaces the standard framework:

\begin{table}[ht]
\centering
\caption{Comparison of standard QM and the crystallisation framework.}
\label{tab:comparison}
\begin{tabular}{@{}ll@{}}
\toprule
Standard QM & Crystallisation framework \\
\midrule
Collapse as axiom & Collapse as emergent convergence \\
Probability treated as fundamental at the operational level & Probability potentially epistemic in this framework \\
Measurement as primitive & Measurement as constraint interaction \\
State selection unspecified & State selection by extremisation of $F$ \\
Indeterminism standardly treated as irreducible & Determinism at the selection level (given $F$) \\
\bottomrule
\end{tabular}
\end{table}

The Dirac equation remains the governing equation for admissible states. The crystallisation architecture adds a selection layer; it does not modify the dynamics of $\mathcal{S}$.

Whether the Born rule can be recovered from ignorance over constraint-landscape variables remains open. At the effective level, the Born rule is retained operationally: if the observer cannot evaluate $F$, the squared-amplitude distribution remains the best available prediction. An epistemic reconstruction of Born statistics from the crystallisation framework is conjectured but not derived in the present paper.

%==================================================================
\section{Geometric Structure}\label{sec:geometry}
%==================================================================

The constraint functional $F$ defines a landscape over $\mathcal{S}_{\mathrm{adm}}$. The topology of this landscape determines the structure of physically realised states:

\begin{itemize}
    \item \textbf{Constraint manifold:} the subspace of $\mathcal{S}_{\mathrm{adm}}$ where $D(\psi) = 0$ (all constraints exactly satisfied). Physically realised states lie on or near this manifold.
    \item \textbf{Attractor basins:} regions of $\mathcal{S}_{\mathrm{adm}}$ that flow toward a specific local maximum of $F$ under the dynamics~\eqref{eq:flow}. In the intended interpretation, distinct basins would correspond to distinct stable realised state classes, potentially including particle types.
    \item \textbf{Basin boundaries:} separatrices between attractor basins. Transitions between basins would correspond to state changes in the physical system.
\end{itemize}

A deeper geometric analysis --- connecting $F$ to specific symmetry structures such as polytope geometry --- is developed in companion work~\citep{SmartIV, SmartV} but is beyond the scope of the present paper.

%==================================================================
\section{Testable Implications}\label{sec:testable}
%==================================================================

The crystallisation architecture generates several structural implications that distinguish it from axiomatic collapse and stochastic-collapse models. These are conceptual implications of the framework; their operationalisation requires further development of $F$:

\begin{enumerate}
    \item \textbf{Preferred states.} Physically realised states correspond to extrema of a constraint functional. If $F$ can be characterised experimentally, preferred states should exhibit higher coherence measures than generic superpositions.

    \item \textbf{Stability spectra.} The set of candidate stable particles corresponds to local maxima of $F$. The stability spectrum is linked to the constraint geometry rather than to arbitrary parameter choices.

    \item \textbf{Deterministic selection signatures.} If the full constraint state were exactly replicated, the framework predicts reproducible selection. In practical experiments, apparent statistical spread may reflect uncontrolled variation in the effective constraint environment.

    \item \textbf{Transition-time structure.} A dynamical implementation of crystallisation via Eq.~\eqref{eq:flow} would naturally introduce a finite convergence timescale, governed by the gradient structure of $F$.

    \item \textbf{Connection to mass structure.} If the constraint functional $F$ is related to the spectral geometry of the underlying space, then the mass ratios of stable particles should be computable from the constraint landscape. Companion work~\citep{SmartIV, SmartV} explores this direction.
\end{enumerate}

%==================================================================
\section{Limitations}\label{sec:limits}
%==================================================================

\begin{itemize}
    \item \textbf{The functional form of $F$ is not uniquely derived.} The constraint functional is a schematic template, not derived from a unique principle. Multiple functionals could produce similar selection behaviour.

    \item \textbf{The coefficients $a_1, a_2, a_3, a_4$ are not determined.} The relative weighting of the functional components is a structural choice within the framework.

    \item \textbf{Lorentz covariance of $F$ has not been established.} Whether the schematic functional can be formulated in a manifestly covariant way is an open question.

    \item \textbf{Compatibility with QFT is not fully established.} The present paper works at the level of the Dirac equation. Extension to the full QFT formalism (operator algebras, renormalisation, gauge invariance) remains to be developed.

    \item \textbf{The Born rule recovery is conjectured, not proven.} Whether Born-rule statistics can be reconstructed from ignorance over constraint-landscape variables is an open mathematical question.

    \item \textbf{Existence and uniqueness of the flow limit are not proven.} The convergence of the flow~\eqref{eq:flow} depends on the specific form of $F$ and the topology of $\mathcal{S}_{\mathrm{adm}}$.

    \item \textbf{The ontology of the evolving state is not specified beyond the formal level.} What physical entity undergoes the flow~\eqref{eq:flow} is an interpretive question not resolved here.

    \item \textbf{This is not a replacement for the Standard Model.} The crystallisation framework addresses the selection mechanism, not the dynamical content of the Standard Model. It is an additional layer, not a substitute.
\end{itemize}

%==================================================================
\section{Conclusion}\label{sec:conclusion}
%==================================================================

The Dirac equation defines the space of admissible relativistic quantum states. Standard quantum mechanics assigns probabilities to outcomes within this space but does not provide a universally accepted physical mechanism for the selection of a definite state. We have proposed crystallisation --- a deterministic, constraint-based selection architecture --- as a candidate resolution.

The architecture defines a projected gradient flow~\eqref{eq:flow} on a restricted admissible domain $\mathcal{S}_{\mathrm{adm}}$, driven by a schematic constraint functional $F$. Within the intended interpretation: candidate stable particles correspond to local maxima of $F$; antiparticles are symmetry-related extrema on conjugate solution branches; and annihilation is heuristically interpreted as a transition toward a more stable configuration in a field-theoretic extension.

The proposal extends rather than replaces the Dirac equation and standard quantum mechanics. The constraint functional is schematic and not uniquely derived; the Lorentz covariance of $F$, the convergence of the flow, and the Born-rule recovery are open questions. The paper introduces a \textit{selection architecture} --- the formal structure within which these questions can be posed --- not a finished theory.

%==================================================================
\appendix
\section{Comparison to Collapse Interpretations}\label{app:comparison}
%==================================================================

\begin{table}[ht]
\centering
\caption{Crystallisation compared to existing interpretations.}
\label{tab:interpretations}
\begin{tabular}{@{}llll@{}}
\toprule
Feature & Copenhagen & GRW & Crystallisation \\
\midrule
Selection mechanism & Axiomatic & Stochastic & Constraint-based ($\arg\sup F$) \\
Randomness & Operational & Fundamental & Epistemic (conjectured) \\
Collapse dynamics & Unspecified & Finite rate & Gradient flow on $F$ \\
Preferred basis & Unspecified & Position & Extrema of $F$ \\
Testable predictions & No (interpretive) & Yes (rate) & Yes (coherence, timescale) \\
Extends Dirac & Adds axiom & Adds noise & Adds constraint layer \\
\bottomrule
\end{tabular}
\end{table}

The key distinction is that crystallisation proposes a deterministic selection architecture grounded in a constraint functional, whereas Copenhagen provides operational closure without mechanism and GRW provides a stochastic modification. All three are compatible with the Dirac equation as the governing equation for admissible states.

%==================================================================
\bibliographystyle{plainnat}
\bibliography{references}

\end{document}
